\documentclass[10pt,a4paper]{amsart}
\usepackage{amsmath,amssymb,amsthm,graphicx}
% Retained minimal preamble for the Stacks Project chapter more-morphisms.tex (More on Morphisms).
% Verbatim excerpt of the chapter preamble: Xy-pic setup (native preamble lines 14-18) and the native
% theorem declarations with the mathematical macros (native preamble lines 154-217). Copied mechanically;
% nothing added, removed or redefined. The wrapper class and the amsmath/amssymb/amsthm/graphicx base
% packages are supplied by the assembler, so the class/externaldocument part of the native preamble is omitted.
% For commutative diagrams we use Xy-pic
\usepackage[all]{xy}

% We use 2cell for 2-commutative diagrams.
\xyoption{2cell}
% Theorem environments.
%
\theoremstyle{plain}
\newtheorem{theorem}[subsection]{Theorem}
\newtheorem{proposition}[subsection]{Proposition}
\newtheorem{lemma}[subsection]{Lemma}

\theoremstyle{definition}
\newtheorem{definition}[subsection]{Definition}
\newtheorem{example}[subsection]{Example}
\newtheorem{exercise}[subsection]{Exercise}
\newtheorem{situation}[subsection]{Situation}

\theoremstyle{remark}
\newtheorem{remark}[subsection]{Remark}
\newtheorem{remarks}[subsection]{Remarks}

\numberwithin{equation}{subsection}

% Macros
%
\def\lim{\mathop{\mathrm{lim}}\nolimits}
\def\colim{\mathop{\mathrm{colim}}\nolimits}
\def\Spec{\mathop{\mathrm{Spec}}}
\def\Hom{\mathop{\mathrm{Hom}}\nolimits}
\def\Ext{\mathop{\mathrm{Ext}}\nolimits}
\def\SheafHom{\mathop{\mathcal{H}\!\mathit{om}}\nolimits}
\def\SheafExt{\mathop{\mathcal{E}\!\mathit{xt}}\nolimits}
\def\Sch{\mathit{Sch}}
\def\Mor{\mathop{\mathrm{Mor}}\nolimits}
\def\Ob{\mathop{\mathrm{Ob}}\nolimits}
\def\Sh{\mathop{\mathit{Sh}}\nolimits}
\def\NL{\mathop{N\!L}\nolimits}
\def\CH{\mathop{\mathrm{CH}}\nolimits}
\def\proetale{{pro\text{-}\acute{e}tale}}
\def\etale{{\acute{e}tale}}
\def\QCoh{\mathit{QCoh}}
\def\Ker{\mathop{\mathrm{Ker}}}
\def\Im{\mathop{\mathrm{Im}}}
\def\Coker{\mathop{\mathrm{Coker}}}
\def\Coim{\mathop{\mathrm{Coim}}}

% Boxtimes
%
\DeclareMathSymbol{\boxtimes}{\mathbin}{AMSa}{"02}

%
% Macros for moduli stacks/spaces
%
\def\QCohstack{\mathcal{QC}\!\mathit{oh}}
\def\Cohstack{\mathcal{C}\!\mathit{oh}}
\def\Spacesstack{\mathcal{S}\!\mathit{paces}}
\def\Quotfunctor{\mathrm{Quot}}
\def\Hilbfunctor{\mathrm{Hilb}}
\def\Curvesstack{\mathcal{C}\!\mathit{urves}}
\def\Polarizedstack{\mathcal{P}\!\mathit{olarized}}
\def\Complexesstack{\mathcal{C}\!\mathit{omplexes}}
% \Pic is the operator that assigns to X its picard group, usage \Pic(X)
% \Picardstack_{X/B} denotes the Picard stack of X over B
% \Picardfunctor_{X/B} denotes the Picard functor of X over B
\def\Pic{\mathop{\mathrm{Pic}}\nolimits}
\def\Picardstack{\mathcal{P}\!\mathit{ic}}
\def\Picardfunctor{\mathrm{Pic}}
\def\Deformationcategory{\mathcal{D}\!\mathit{ef}}
\newenvironment{reference}{}{}
\newenvironment{slogan}{}{}
\newenvironment{history}{}{}
\usepackage{cleveref}
\crefname{originallemma}{original lemma}{original lemmas}
\Crefname{originallemma}{Original lemma}{Original lemmas}
\makeatletter
\newlabel{algebra-lemma-isolated-point-fibre}{{00PK}{source}{}{sourcepointer-1}{}}
\newlabel{algebra-lemma-isolated-point-fibre@cref}{{[originallemma][0][]00PK}{[1][0][]source}{}{}{}}
\makeatother
\crefname{originallemma}{original lemma}{original lemmas}
\Crefname{originallemma}{Original lemma}{Original lemmas}
\makeatletter
\newlabel{algebra-lemma-quasi-finite-open}{{00QA}{source}{}{sourcepointer-2}{}}
\newlabel{algebra-lemma-quasi-finite-open@cref}{{[originallemma][0][]00QA}{[1][0][]source}{}{}{}}
\makeatother
\crefname{originallemma}{original lemma}{original lemmas}
\Crefname{originallemma}{Original lemma}{Original lemmas}
\makeatletter
\newlabel{properties-lemma-complement-closed-point-Jacobson}{{02IM}{source}{}{sourcepointer-3}{}}
\newlabel{properties-lemma-complement-closed-point-Jacobson@cref}{{[originallemma][0][]02IM}{[1][0][]source}{}{}{}}
\makeatother
\title{A local ring map with finite residue extension admits a prime outside the closed fibre whose quotient is quasi-finite over the base}
\author{The Stacks Project Authors}
\date{stacks-project:more-morphisms.tex:local-snapshot-20261011; GNU FDL 1.2 or later}
\begin{document}
\maketitle
\noindent\textit{Editorial note.} Statement span is lines 14957-14967 of sources/stacks-more-morphisms/pinned-original-source.tex; proof span is lines 14969-14991, the complete author proof (21 body lines).\par\medskip
\noindent\textit{Editorial note.} Navigation stubs: 3 cited-but-not-included labels are supplied as non-load-bearing pointers with their pinned tags where the snapshot tag index has them.\par\medskip
\noindent\textit{Editorial note.} extra\_packages is empty on purpose: the retained preamble already loads xy with options and a second declaration clashes in the assembler wrapper.\par\medskip
\noindent\textit{Editorial note.} Text is retained verbatim, including the author's own typos and placeholders.\par\medskip
\noindent\textit{Editorial note.} First of the chapter's quasi-section pair; the sibling card (nearby-open-X variant) sits further down the queue and is a different statement.\par\medskip
\noindent\textit{Allowed standard prerequisites (editorial).}\begin{itemize}
\item Explicit prerequisite boundary. Jacobian-ness of the punctured spectrum of a local Noetherian ring (Properties, lemma-complement-closed-point-Jacobson), the isolated-point fibre criterion for quasi-finiteness (Algebra, lemma-isolated-point-fibre), and openness of the quasi-finite locus (Algebra, lemma-quasi-finite-open).
\end{itemize}\medskip
\ifdefined\hypertarget\hypertarget{sourcepointer-1}{}\fi
\noindent\textit{Source reference.} Original lemma 00PK: sources/stacks-more-morphisms/pinned-original-source.tex, source lines 14991--14991. Cited by the retained statement or proof but not part of the collected spans; supplied only as a navigation pointer. The argument does not depend on its content.\par\medskip
\ifdefined\hypertarget\hypertarget{sourcepointer-2}{}\fi
\noindent\textit{Source reference.} Original lemma 00QA: sources/stacks-more-morphisms/pinned-original-source.tex, source lines 14991--14991. Cited by the retained statement or proof but not part of the collected spans; supplied only as a navigation pointer. The argument does not depend on its content.\par\medskip
\ifdefined\hypertarget\hypertarget{sourcepointer-3}{}\fi
\noindent\textit{Source reference.} Original lemma 02IM: sources/stacks-more-morphisms/pinned-original-source.tex, source lines 14991--14991. Cited by the retained statement or proof but not part of the collected spans; supplied only as a navigation pointer. The argument does not depend on its content.\par\medskip
\begin{lemma}
\label{lemma-quasi-finite-quasi-section-meeting-nearby-open}
Let $\varphi : A \to B$ be a local ring map of local rings.
Let $V \subset \Spec(B)$ be an open subscheme
which contains at least one prime not lying over $\mathfrak m_A$.
Assume $A$ is Noetherian, $\varphi$ essentially of finite type, and
$A/\mathfrak m_A \subset B/\mathfrak m_B$ is finite.
Then there exists a $\mathfrak q \in V$,
$\mathfrak m_A \not = \mathfrak q \cap A$ such that
$A \to B/\mathfrak q$ is the localization of a quasi-finite ring map.
\end{lemma}

\begin{proof}
Since $A$ is Noetherian and $A \to B$ is essentially of finite type,
we know that $B$ is Noetherian too. By
Properties, Lemma \ref{properties-lemma-complement-closed-point-Jacobson}
the topological space $\Spec(B) \setminus \{\mathfrak m_B\}$
is Jacobson. Hence we can choose a closed point $\mathfrak q$
which is contained in the nonempty open
$$
V \setminus \{\mathfrak q \subset B \mid \mathfrak m_A = \mathfrak q \cap A\}.
$$
(Nonempty by assumption, open because $\{\mathfrak m_A\}$ is a closed
subset of $\Spec(A)$.)
Then $\Spec(B/\mathfrak q)$ has two points, namely $\mathfrak m_B$
and $\mathfrak q$ and $\mathfrak q$ does not lie over $\mathfrak m_A$.
Write $B/\mathfrak q = C_{\mathfrak m}$ for some finite type $A$-algebra
$C$ and prime ideal $\mathfrak m$. Then $A \to C$ is quasi-finite at
$\mathfrak m$ by
Algebra, Lemma \ref{algebra-lemma-isolated-point-fibre} (2).
Hence by
Algebra, Lemma \ref{algebra-lemma-quasi-finite-open}
we see that after replacing $C$ by a principal localization the ring
map $A \to C$ is quasi-finite.
\end{proof}

\end{document}
